\chapter{Dialling in settings}
\label{ch:tuning}

Every laser, every tube, every diode and every sheet of material is slightly
different. The recipes you find on the internet are starting points for someone
else's machine. This chapter is about replacing them with numbers that belong to
yours --- systematically, and in an afternoon rather than over a year of
frustration.

\section{The method}

Three rules make all the difference:

\begin{enumerate}[leftmargin=1.6em]
  \item \textbf{Change one thing at a time.} A grid with speed on one axis and
        power on the other, all burned in a single job on a single piece of
        material, tells you more than fifty separate experiments.
  \item \textbf{Write down what you burned.} The grid is worthless without a
        record of the settings and the result. Engrave the labels into the test
        piece, or photograph it with the numbers beside it.
  \item \textbf{Test the material you will actually use.} Thickness, moisture,
        glue content and the colour of the sheet all change the answer. Plywood
        from two suppliers is two materials.
\end{enumerate}

\section{Test grids: the Parameter-Test window}

\menu{Tools>Laser-Tools>Parameter-Test} builds the grid for you. It has two tabs:
\tabname{Generator}, where the grid is designed, and \tabname{Templates}, where
you save and reload designs.

\begin{walkthrough}{Burn a speed/power grid on 3\unit{mm} plywood}
\begin{enumerate}[label=\arabic*.,leftmargin=1.4em]
  \item Open \menu{Tools>Laser-Tools>Parameter-Test}.
  \item In \emph{Operation to test}, choose \btn{Cut}, \btn{Engrave},
        \btn{Raster}, \btn{Image}, \btn{Hatch} or \btn{Wobble}. The parameter
        rows below change with the choice: Cut and Engrave give you
        \emph{Speed}, \emph{Power} and \emph{Passes}; Raster and Image add
        \emph{DPI} and \emph{Overscan}; Hatch adds \emph{Hatch Distance} and
        \emph{Hatch Angle}; Wobble adds its four wobble parameters.
  \item Fill in \emph{First parameter (X-Axis)} and \emph{Second parameter
        (Y-Axis)}: which parameter, how many steps (\btn{Count}), the
        \btn{Minimum} and \btn{Maximum}, the cell \btn{Width}/\btn{Height},
        the \btn{Delta} between steps and the colour. A sensible first grid for
        cutting 3\unit{mm} plywood is speed 5--20\unit{mm/s} in five steps on X,
        and power 40--100\% in five steps on Y.
  \item Decide about labels. \btn{Labels} prints a descriptive label down the
        sides of the grid; \btn{Values} prints the numbers; \btn{Simple} uses
        vector shapes for the labels. Leave both on for your first grid.
  \item Optionally tick \btn{Create boundary shape} to repeat the shape around
        the effect, which makes the grid easier to cut free afterwards.
  \item Press \btn{Create Pattern}. \emph{This deletes all existing operations
        and elements} --- the dialog says so, and asks whether you want to empty
        the project or keep it. On a fresh project, empty it.
  \item Inspect the result on screen: it is a real, editable document, so you can
        move it onto your material's position, or change one row before burning.
  \item Press \btn{Queue to Spooler}. It plans the job in the background and puts
        it in the \panel{Spooler}. It does \emph{not} start it; press
        \btn{Start} in the spooler when the material is in place.
\end{enumerate}
\end{walkthrough}

\begin{figure}[htbp]
\centering
\begin{tikzpicture}[scale=0.92]
\foreach \x in {0,...,4} {
  \foreach \y in {0,...,4} {
    \pgfmathsetmacro{\p}{50+10*\y}
    \pgfmathsetmacro{\s}{5+3*\x}
    \draw[fill=black!\p!white,draw=mkgrey!60] (\x*1.25,\y*1.25) rectangle ++(1.05,1.05);
  }
}
\foreach \x in {0,...,4} {
  \pgfmathsetmacro{\s}{5+3*\x}
  \node[font=\sffamily\tiny,color=mkdark] at (\x*1.25+0.52,-0.28) {\s};
}
\foreach \y in {0,...,4} {
  \pgfmathsetmacro{\p}{50+10*\y}
  \node[font=\sffamily\tiny,color=mkdark] at (-0.42,\y*1.25+0.52) {\p};
}
\node[font=\sffamily\scriptsize,color=mkdark] at (2.6,-0.85) {speed (mm/s) $\rightarrow$};
\node[font=\sffamily\scriptsize,color=mkdark,rotate=90,anchor=south] at (-1.15,2.6) {power (\%) $\rightarrow$};
\end{tikzpicture}
\caption{A completed 5\,$\times$\,5 cut grid. Read it once: the squares that fell
out define your cutting window; the top-right corner is burned through or charred;
the bottom-left does not cut at all. The darkest squares that still cut cleanly
are the ones worth writing down.}
\label{fig:testgrid}
\end{figure}

\subsection*{Reading the grid}

Look for three boundaries rather than for a single best square:

\begin{itemize}
  \item the \textbf{does-not-cut} boundary --- everything below and right of it
        wastes material;
  \item the \textbf{clean-cut} band --- squares that separate and have brown, not
        black, edges;
  \item the \textbf{charring} boundary --- where the edge is black, sticky, or
        the material has a wide heat-affected zone around it.
\end{itemize}

\noindent Pick the setting in the clean-cut band that is nearest to the
does-not-cut boundary: it is the fastest one that works, and it is the least
likely to cause a fire. For engraving, the equivalent goal is the lightest burn
that is consistently dark.

\begin{tipbox}[Save the grid as a template]
In the \tabname{Templates} tab, type a name and press \btn{Save}. You then have
a one-click grid for the next new material, sized the way you like it. Reloading
a template and pressing \btn{Create Pattern} takes ten seconds.
\end{tipbox}

\section{Storing the answer: the Material Library}

A grid gives you numbers; the \btn{Material Library} keeps them.
Open it with the \btn{Material Library} button or \menu{Tools>Material Library}.
An entry is a named collection of operation presets, keyed by
\emph{Material}, \emph{Thickness} and \emph{Laser type}, with an optional note.

\begin{walkthrough}{Turn your grid into a library entry}
\begin{enumerate}[label=\arabic*.,leftmargin=1.4em]
  \item Build the operations you want to keep --- for example a \mkseries{Cut}
        at your measured speed, power and passes, and an \mkseries{Engrave} at
        its own numbers. If you have a saved project that already contains them,
        load it.
  \item Open the Material Library and press \btn{Get current} to capture the
        operations from the tree as a new entry. (Alternatively,
        \btn{Add new} creates an empty entry and you add operations by name from
        the right-click menu: \emph{Add Cut}, \emph{Add Engrave},
        \emph{Add Raster}, \emph{Add Image}, \emph{Add diagonal Hatch},
        \emph{Add Wobble}, and so on.)
  \item Fill in the \emph{Information} box: a good \emph{Title}, the
        \emph{Material} and its \emph{Thickness}, the \emph{Laser} type, the
        laser's \emph{Power} in watts and — for galvo machines — the
        \emph{Lens-Size}. Add anything you would otherwise forget in the note
        field: supplier, tape or no tape, focus setting, whether it needed two
        passes.
  \item Use the filter boxes at the top (\emph{Material}, \emph{Thickness},
        \emph{Laser}, \btn{Reset Filter}) to check your entry lands where you
        expect, and press \btn{Use for statusbar} if this is the entry you want
        the quick-assign squares to offer.
  \item Later: select the entry and press \btn{Load into Tree} to put its
        operations straight into your job.
\end{enumerate}
\end{walkthrough}

\noindent The library imports other people's data too: \btn{Import} reads
EZcad, LightBurn and MeerK40t library files (\file{*.lib}, \file{*.ini},
\file{*.clb}, \file{*.cfg}), with options to compensate for a different lens size
or a different laser power. That is worth knowing when you inherit a machine: the
settings are probably already in a file somewhere, in units and geometry you do
not share, which is exactly what the compensation fields are for.

\begin{notebox}[Keep it honest]
A library is only as good as its entries. If you change the tube, clean the lens,
or switch suppliers, either re-test or mark the affected entries as stale. A
wrong entry is more dangerous than no entry, because it saves time --- right up
until the job fails.
\end{notebox}

\section{Kerf: cutting parts that fit}

The beam removes material. That removed width is the \emph{kerf}, and it is why a
10\unit{mm} tab does not fit a 10\unit{mm} slot.

\subsection*{Measuring your kerf}

\begin{enumerate}[label=\arabic*.,leftmargin=1.4em]
  \item Cut a simple shape of known size --- a 20\unit{mm} square in 3\unit{mm}
        plywood is ideal.
  \item Measure the finished part. If it is 19.85\unit{mm} wide, the beam took
        0.15\unit{mm}: 0.075\unit{mm} from each side.
  \item Measure it in both directions and on two different materials if you cut
        both. Kerf varies with material, focus, speed and power --- and with how
        well the machine tracks.
\end{enumerate}

\noindent Then enter the total width in the operation's
\btn{Kerf compensation} field, for example \cmd{0.15mm}. MeerK40t offsets each
closed path by half that amount so the finished part comes out at the designed
size. For installation instructions on the faster offsetting library, see
Chapter~\ref{ch:install}: without \cmd{pyclipr} installed, offsets are computed
by a slower fallback.

\begin{itemize}
  \item For \emph{fits that must be tight or loose}, do not rely on kerf
        compensation alone: design the clearance into the artwork. A
        ``10\unit{mm}'' slot intended for a 10\unit{mm} tab is always wrong; make
        it 10.1\unit{mm} or 10.2\unit{mm} and cut a test coupon.
  \item For one-off offsets, the element context menu is quicker:
        \menu{Outline element(s)} produces an outline path at a given distance,
        and the \menu{Offset shapes} submenu offsets a shape to the outside or
        the inside --- useful for making a lid that matches a box, or a mask that
        surrounds a part.
  \item Compensation applies per operation, so cut and engrave can differ.
\end{itemize}

\section{The machine's own calibration}

Some wrong numbers are not material settings at all, but the machine's idea of
distance. Calibrate these once, correctly, and they stay correct.

\begin{center}\small
\begin{tabular}{@{}L{3.4cm}L{5.0cm}L{5.7cm}@{}}
\toprule
\textbf{What} & \textbf{Symptom} & \textbf{How to fix} \\
\midrule
Axis scale (\setting{scale_x}, \setting{scale_y}) & A 100\unit{mm} move measures
103\unit{mm} & Jog exactly 100\unit{mm}, measure, and correct the scale factor in
the device configuration by the ratio measured $\div$ commanded. \\
Backlash & Squares are not square, or lines repeat shifted in one direction &
Newly exposes \emph{Horizontal/Vertical backlash} in millimetres; other drivers
need mechanical attention (belt tension, eccentric nuts). \\
Speed chart (Newly) & Corners rounded or overshooting at speed & The device's
\emph{Raster Chart} holds \emph{Speed <=}, \emph{Acceleration Length},
\emph{Backlash} and \emph{Corner Speed} rows. \\
Acceleration (K40/Lihuiyu) & Raster lines ragged or the machine skipping at high
speed & The \btn{Acceleration-Chart} window (\menu{Tools>Acceleration-Chart}) has
separate breakpoints for \emph{Vector}, \emph{Vertical Raster} and
\emph{Horizontal Raster}. \\
Rotary steps (GRBL) & Wrapped artwork stretched or squashed along the
circumference & \cmd{rotarycal <measured_mm>} recalculates steps/mm from a test
burn; \cmd{rotarysuggest <diameter>} estimates it from your mechanics. \\
Home corner & The head homes to the corner you do not expect, and jobs land
mirrored & \emph{Force Declared Home} in the device configuration. \\
\bottomrule
\end{tabular}
\end{center}

\section{Beyond speed and power}

Four settings change results in ways that surprise people.

\begin{description}[leftmargin=1.4em,style=nextline]
  \item[PPI / frequency] Pulses per inch, or kHz on newer machines. Lower PPI
        leaves a visible serration on acrylic edges and gives wood a rougher,
        more charred cut; higher PPI gives smoother edges but can overheat the
        material. If your machine exposes \btn{Frequency}, run one row of your
        test grid across the values rather than guessing.
  \item[Focus and material height] Focus is depth of field, and half of
        ``my settings stopped working'' is a job run a millimetre out of focus.
        Focus on the surface of the material you cut, not on a fixed table
        height, and re-check it between pieces of different thickness.
  \item[Masking and surface protection] Masking tape or transfer film reduces
        scorch marks around a raster and matters more than most settings for
        the final appearance on light-coloured wood and on painted surfaces.
        Remember that tape changes what the beam sees first.
  \item[Air assist] Clearing the smoke changes the burn as much as 10\% of power:
        too little and the cut re-deposits soot and catches fire; too much and a
        delicate engraving is blown pale.
\end{description}

\begin{machinebox}{Keep a record that ages well}
The most valuable page in a workshop is not a settings table; it is a
\emph{sourced} settings table. For each material write down: supplier and
product, thickness measured with callipers, date, machine, focus setting, air
assist on or off, tape or no tape, speed, power, passes, PPI and the verdict in
one word. Six months later those columns are the difference between a repeatable
process and superstition.
\end{machinebox}

\begin{tryit}{Build your own starting kit}
\begin{enumerate}[label=\arabic*.,leftmargin=1.4em]
  \item Pick the material you use most. Burn a 5$\times$5 cut grid and a
        3$\times$3 engrave grid on one piece.
  \item Photograph both with a written card in shot, and keep the photographs
        with the physical samples in a labelled bag.
  \item Measure the kerf from the cut grid and record it.
  \item Enter the winning settings into the Material Library as two entries: one
        for cutting, one for engraving.
  \item Repeat for your second material. Then stop: two well-characterised
        materials beat ten guesses.
  \item From now on, when you have to guess on a new material, guess by
        interpolation from these two --- and test the guess on scrap before it
        touches a real job.
\end{enumerate}
\end{tryit}
